\documentclass[12pt]{article}
\usepackage[utf8]{inputenc}
\usepackage[T1]{fontenc}
\usepackage[margin=1in]{geometry}
\usepackage{setspace}
\usepackage{xcolor}
\usepackage[colorlinks=true, urlcolor=blue!60!black, linkcolor=blue!60!black]{hyperref}
\setlength{\parindent}{0em}
\setlength{\parskip}{0.45em}
\renewcommand{\baselinestretch}{1.0}
\begin{document}
\begin{center}
{\Large \textit{Reed v. Town of Gilbert}, 576 U.S. 155 (2015)\par}
\vspace{0.3em}
{\normalsize Case Brief --- One-Page Edition\par}
\end{center}

\textbf{Memory Jogger.} The church-sign case: facial subject-matter distinctions trigger strict scrutiny (\href{https://supreme.justia.com/cases/federal/us/576/155/}{576 U.S. 155}).

\textbf{Facts.} Gilbert, Arizona let ideological and political signs be larger and stay up far longer than temporary directional signs for religious or nonprofit meetings, which were limited to 6 square feet and 12 hours before the event. Pastor Clyde Reed's church was cited for its Sunday-service signs.

\textbf{Procedural History.} The Ninth Circuit twice upheld the Code as content neutral for lack of censorial motive.

\textbf{Votes.} Reversed, 9--0 in the judgment; Thomas wrote for six; Alito concurred; Breyer and Kagan concurred in the judgment.

\textbf{Holding.} A law that distinguishes speech on its face by topic, function, or purpose is content based and subject to strict scrutiny regardless of benign motive; Gilbert's underinclusive Code failed it.

\textbf{Reasoning.} A sign's treatment turned entirely on its message; motive matters only for facially neutral laws, and subject-matter discrimination is barred even without viewpoint discrimination. Alito listed neutral rules towns may keep; Kagan warned that courts would become a ``veritable Supreme Board of Sign Review.''

\textbf{Examples.} \emph{Same-Side:} a ban exempting only real estate and garage-sale signs. \emph{Opposite-Side:} uniform size limits and off-premises billboard rules. \emph{Fence-Sitter:} a uniform limit on one-time event signs.

\textbf{Critique.} Strict formalism imperils benign safety rules, and Alito's concurrence sat in tension with Thomas's reasoning until \textit{City of Austin}.

\textbf{Key Quotation.} ``Content-based laws---those that target speech based on its communicative content---are presumptively unconstitutional and may be justified only if the government proves that they are narrowly tailored to serve compelling state interests.'' --- Justice Thomas, Opinion of the Court, 576 U.S.\ at 163.

\end{document}
